\section{Evaluasi Metodologi \& Koreksi Matematis}

\begin{frame}{Reformulasi Hamiltonian \& Lagrangian}
    \begin{itemize}
        \item \textbf{Koreksi Tanda Hamiltonian:} Persamaan Hamiltonian disesuaikan menjadi negatif untuk merepresentasikan minimalisasi energi yang konsisten dengan formulasi Ising Model.
        \item \textbf{Eliminasi Double Counting:} Perbaikan pada perhitungan matriks kovarians $\Sigma_{ij}$ untuk memastikan tidak ada penghitungan ganda pada interaksi antar aset.
        \item \textbf{Lagrangian Multiplier Dinamis ($A$):}
        \begin{equation}
            A(N) = 2.0 + 2 \ln(1 + N)
        \end{equation}
        \item \textbf{Tujuan Nilai $A$ Dinamis:}
        \begin{itemize}
            \item \textbf{Stabilisasi Lanskap Energi:} Menjaga rasio antara suku \textit{return-risk} dan suku penalti agar tetap seimbang seiring bertambahnya dimensi $N$.
            \item \textbf{Optimalisasi Constraint:} Mencegah suku penalti mendominasi Hamiltonian secara berlebihan (klasik pada $A$ statis) yang dapat menyebabkan VQE terjebak dalam solusi trivial (semua nol).
        \end{itemize}
    \end{itemize}
\end{frame}

\begin{frame}{Justifikasi Hukum Logaritmik Penalti $A(N)$}
    \begin{itemize}
        \item \textbf{Skalabilitas Ruang Konfigurasi:} Ruang pencarian Ising Model tumbuh secara eksponensial $2^N$. Penalti statis gagal menjaga integritas solusi pada $N$ besar.
        \item \textbf{Analogi Entropi Konfigurasi:}
        \begin{itemize}
            \item Berdasarkan prinsip \textit{Maximum Entropy}, ketidakpastian sistem tumbuh secara logaritmik terhadap jumlah kemungkinan \textit{state}.
            \item Penggunaan $\ln(1+N)$ menyelaraskan kekuatan \textit{constraint} dengan laju pertumbuhan kepadatan energi dalam sistem kuantum diskret.
        \end{itemize}
        \item \textbf{Keseimbangan Eksplorasi-Eksploitasi:} Skala logaritmik memastikan kenaikan penalti cukup kuat untuk mematuhi batasan portofolio, namun cukup halus untuk tidak menghancurkan gradien yang diperlukan untuk konvergensi ke \textit{Ground State}.
    \end{itemize}
\end{frame}

\begin{frame}{Motivasi Penggunaan Entanglement}
    \begin{itemize}
        \item \textbf{Korelasi Aset Sebagai Entanglement:} Dalam konteks ekonomi fisik, korelasi antar aset tidak bersifat independen. Entanglement digunakan untuk memodelkan hubungan non-lokal antar variabel finansial.
        \item \textbf{Representasi State Kuantum:} 
        \begin{itemize}
            \item Memungkinkan eksplorasi ruang konfigurasi yang lebih luas dibandingkan metode stokastik klasik.
            \item Menangkap dependensi probabilitas antar qubit (aset) yang merepresentasikan keputusan investasi kolektif.
        \end{itemize}
        \item \textit{"Tanpa entanglement, sirkuit VQE hanya akan merepresentasikan distribusi probabilitas marginal yang independen."}
    \end{itemize}
\end{frame}
